\subsection{Two Comings, and the One Between}\label{sec:comings}

The first Father the Church reads in Advent is St Cyril, later bishop of
Jerusalem, teaching beginners. He gave his lectures to men and women preparing
for Baptism, and on the First Sunday the Office of Readings reads one of them:

\begin{quote}
We preach not one advent only of Christ, but a second also, far more glorious
than the former. For the former gave a view of His patience; but the latter
brings with it the crown of a divine kingdom. For all things, for the most part,
are twofold in our Lord Jesus Christ: a twofold generation; one, of God, before
the ages; and one, of a Virgin, at the close of the ages: His descents twofold;
one, the unobserved, like rain on a fleece; and a second His open coming, which
is to be. In His former advent, He was wrapped in swaddling clothes in the
manger; in His second, He covereth Himself with light as with a garment. In His
first coming, He endured the Cross, despising shame; in His second, He comes
attended by a host of Angels, receiving glory. We rest not then upon His first
advent only, but look also for His second.\par
{\smallskip\noindent\small ---St Cyril of Jerusalem, Catechesis 15\par}
\end{quote}

\noindent Cyril goes on to find both comings in the prophet Malachi, and then he
puts it plainly: ``Our Lord Jesus Christ, then, comes from heaven; and He comes
with glory at the end of this world, in the last day.''

The Church still teaches Advent in Cyril's terms. The Universal Norms on the
Liturgical Year, issued with the reformed calendar in 1969, give the season a
twofold character (n.~39). It prepares for Christmas, when the first coming of
the Son of God among men is remembered; and by that remembrance it turns the
mind to his second coming at the end of time. For both reasons it is a season
of devout and joyful expectation. The commentary published with the new
calendar in 1969 adds that Advent kept its four weeks but was no longer to be
held a season of penance only. The Second Vatican Council, when it lists the
mysteries of Christ that the Church unfolds through the year, ends the list with
the hope of his coming; and the coming meant there is the last one
(\emph{Sacrosanctum Concilium}, n.~102).

The Missal marks the two halves of the season with two Prefaces. Until 16
December the priest sings the first, which speaks of both comings: the Lord came
first in humble flesh, and will come again in his majesty. From 17 December the
second Preface turns to Christmas itself, naming the prophets who foretold him,
the Virgin Mother who awaited him with a love no words can tell, and John who
pointed him out when he came.

The first Preface keeps the frame of a far older prayer. In the Veronese
Sacramentary, an ancient collection of Roman Mass prayers, a preface for the Ascension
asked that our minds be raised to where the Redeemer had gone up,
\latin{ut secundo Mediatoris adventu manifesto munere capiamus quod nunc audemus
sperare promissum}, ``so that at the second coming of the Mediator we may
receive as an open gift what we now dare to hope for as promised.'' The Advent
Preface keeps that hope of the second coming and adds the word for waiting: the
faithful now await, watching, what was promised.

Between the two comings St Bernard found a third. His fifth sermon for Advent,
read in the Office on the Wednesday of the first week, speaks of it:

\begin{quote}
This intermediate place is in fact another advent wherein such as have knowledge
thereof enjoy a most pleasant repose. For whereas the other two are known to
all, this is secret. \ldots{} But in that which intervenes He is hidden, visible
only to the elect who see Him in themselves, and so their souls are saved.
Therefore, in the first He came in the flesh and in weakness, in the second He
comes in spirit and in virtue, and in the third He shall come in glory and
majesty. \ldots{} Thus, this intermediate advent is the way, so to speak, by
which we must travel from the first to the last. In the first Christ was our
Redemption, in the third He shall appear as our Life, whilst in this second He
is our Repose and our Consolation.\par
{\smallskip\noindent\small ---St Bernard, Sermon 5 for Advent\par}
\end{quote}

\noindent A little further on in the same sermon he tells his monks how it is
received:

\begin{quote}
My brethren, if you thus keep the word of God, there can be no doubt that it, in
turn, shall keep you. For the Son will come to you with the Father, the great
Prophet will come to you to rebuild the Jerusalem of your souls and to make all
things new.\par
{\smallskip\noindent\small ---St Bernard, Sermon 5 for Advent\par}
\end{quote}

On 21 December, in the last days before Christmas, the Collect of the day
asks that those who are glad of the Lord's coming in the flesh may have eternal
life as their reward when he comes in majesty. The prayer grew from one in the Gelasian Sacramentary, which asked that those who
\latin{de adventu Unigeniti tui secundum carnem laetantur, in secundo, cum venerit
in maiestate sua, praemium aeternae vitae percipiant}, ``rejoice at the coming
of thy Only-begotten according to the flesh, may at the second, when he comes in
his majesty, receive the reward of eternal life.'' Advent keeps both comings in
view to the end, and the middle coming is the one a Christian lives in every
day.
